% ----------------------------------------------------------------------
% Parte I, capítulo 3
% ----------------------------------------------------------------------
\chapter{Dimensões da interoperabilidade}
\label{cap:interoperabilidade}

\instrucao{Este é o capítulo em que o referencial vira modelo. Termine com a proposta original da tese, de preferência em figura, e com as definições operacionais que serão medidas na Parte~\ref{part:empirica}.}

\section{Três dimensões}
A literatura distingue ao menos três dimensões. A sintática trata do formato e do protocolo de troca. A semântica trata do significado dos elementos e dos valores. A organizacional trata de políticas, responsabilidades e acordos entre instituições.

\section{Modelo proposto}
\instrucao{Apresente o modelo com uma figura e uma equação, se houver índice. Defina cada símbolo e remeta à lista de símbolos.}

A Figura~\ref{fig:modelo} mostra o modelo de avaliação proposto. O índice geral combina as três dimensões com pesos definidos por especialistas:
\begin{equation}
  I = \sum_{k \in \{s,m,o\}} w_k\, I_k, \qquad \sum_k w_k = 1,
  \label{eq:indice}
\end{equation}
em que $I_k$ é o índice da dimensão $k$, entre 0 e 1.

\begin{figure}[htb]
  \centering
  \caption{Modelo de avaliação da interoperabilidade de metadados}
  \label{fig:modelo}
  \begin{tikzpicture}[
      dim/.style={draw=vinho, thick, fill=vinho!7, minimum width=3.6cm, minimum height=1.6cm,
        align=center, font=\small},
      seta/.style={-{Stealth}, thick, vinho}]
    \node[dim] (s) {Sintática\\ \footnotesize protocolo e formato};
    \node[dim, right=0.6cm of s] (m) {Semântica\\ \footnotesize vocabulários e autoridades};
    \node[dim, right=0.6cm of m] (o) {Organizacional\\ \footnotesize políticas e acordos};
    \node[draw=vinho, thick, rounded corners, fit=(s)(m)(o), inner sep=8pt] (caixa) {};
    \node[draw=vinho, fill=vinho, text=white, font=\small\bfseries, below=1cm of m,
      minimum width=4cm, minimum height=0.8cm] (I) {Índice geral $I$};
    \draw[seta] (caixa.south) -- (I.north);
  \end{tikzpicture}
  \fonte{Elaborada pelo autor.}
\end{figure}
